% ============================================================
% Capitolo 2 — Background (traduzione italiana di
% chapters/ch02_background.tex; label e citazioni invariati)
% ============================================================
\chapter{Background}
\label{ch:background}

Questo capitolo introduce i concetti su cui il resto della tesi si
costruisce: apprendimento per rinforzo a singolo agente e gradienti di
policy (\cref{sec:rl-essentials}), PPO e stima del vantaggio
(\cref{sec:ppo}), apprendimento multi-agente e MAPPO (\cref{sec:marl}),
curriculum learning e il suo controllo ad addestramento misto
(\cref{sec:curriculum-bg}) e il simulatore CARLA
(\cref{sec:carla-bg}). Ogni concetto è presentato solo alla profondità
alla quale viene usato in seguito.

% ------------------------------------------------------------
\section{Elementi di apprendimento per rinforzo}
\label{sec:rl-essentials}

L'apprendimento per rinforzo (RL) formalizza la decisione sequenziale
come \emph{processo decisionale di Markov} (MDP): una tupla
$(\mathcal{S}, \mathcal{A}, P, r, \gamma)$ di stati, azioni, dinamica di
transizione $P(s' \mid s, a)$, funzione di reward $r(s, a)$ e fattore di
sconto $\gamma \in [0, 1)$~\cite{sutton2018rl}. Un agente agisce tramite
una \emph{policy} $\pi(a \mid s)$, e l'obiettivo di apprendimento è
massimizzare il ritorno atteso scontato
\begin{equation}
  J(\pi) \;=\; \mathbb{E}_{\pi}\!\left[\sum_{t=0}^{\infty} \gamma^{t}\,
  r(s_t, a_t)\right].
\end{equation}
Lo sconto fissa un orizzonte decisionale effettivo di circa
$1/(1-\gamma)$ passi. Non è una costante cosmetica: la piattaforma del
\cref{ch:system} usa $\gamma = 0.997$, cioè un orizzonte effettivo di
$\approx 333$ passi, commisurato a episodi fino a 1{.}000 passi di
controllo in cui reward sparse di progresso sul percorso possono
distare decine di secondi l'una dall'altra.

Due funzioni di valore riassumono una policy: il valore di stato
$V^{\pi}(s) = \mathbb{E}_{\pi}[\,\text{ritorno} \mid s_0 = s\,]$ e il
valore di azione $Q^{\pi}(s, a)$; la loro differenza
$A^{\pi}(s, a) = Q^{\pi}(s, a) - V^{\pi}(s)$, il \emph{vantaggio},
misura quanto un'azione sia migliore della media della policy in quello
stato.

\paragraph{Gradienti di policy e metodi actor--critic.}
Quando la policy è una funzione differenziabile $\pi_{\theta}$ (qui, una
rete neurale), il teorema del gradiente di policy fornisce una stima non
distorta del gradiente dell'obiettivo:
\begin{equation}
  \nabla_{\theta} J \;=\;
  \mathbb{E}_{\pi_{\theta}}\!\left[\nabla_{\theta}
  \log \pi_{\theta}(a_t \mid s_t)\; A^{\pi}(s_t, a_t)\right].
\end{equation}
Stimare il vantaggio con ritorni Monte Carlo grezzi produce alta
varianza; sottrarre una baseline appresa di valore di stato riduce la
varianza senza introdurre bias. I metodi \emph{actor--critic} addestrano
perciò due funzioni: l'\emph{actor} $\pi_{\theta}$ e un \emph{critic}
$V_{\phi}$ usato per costruire le stime di vantaggio. Questa divisione
del lavoro è l'aggancio per l'architettura multi-agente del
\cref{sec:marl}: il critic è un dispositivo del solo addestramento, e
può quindi ricevere informazione che l'actor non vede mai.

\paragraph{Azioni continue.}
Per il controllo continuo la policy è comunemente una gaussiana
diagonale: la rete emette una media (e una log-deviazione standard) per
dimensione d'azione, e le azioni sono campionate attorno a essa. Il
campionamento non è un dettaglio implementativo ma parte dell'oggetto
appreso --- valutare una tale policy sulla sua media deterministica ne
scarta la struttura stocastica, un punto che diventa concreto nella
\cref{sec:bugs}.

% ------------------------------------------------------------
\section{Proximal Policy Optimization}
\label{sec:ppo}

Gli aggiornamenti di gradiente di policy ``vanilla'' sono fragili: un
singolo passo troppo grande può far collassare una policy funzionante, e
i dati on-policy sono obsoleti subito dopo l'aggiornamento.
PPO~\cite{schulman2017ppo} affronta entrambi i problemi con un
\emph{obiettivo surrogato clippato}. Con il rapporto di importanza
$\rho_t(\theta) = \pi_{\theta}(a_t \mid s_t) / \pi_{\theta_{\text{old}}}(a_t \mid s_t)$,
PPO massimizza
\begin{equation}
  L^{\mathrm{CLIP}}(\theta) \;=\;
  \mathbb{E}_t\!\left[\min\bigl(\rho_t A_t,\;
  \operatorname{clip}(\rho_t,\, 1-\varepsilon,\, 1+\varepsilon)\, A_t\bigr)\right],
\end{equation}
che rimuove ogni incentivo a spostare il rapporto oltre
$[1-\varepsilon, 1+\varepsilon]$, permettendo più epoche di passi di
gradiente su minibatch per ogni batch di esperienza senza aggiornamenti
distruttivi. Le implementazioni aggiungono spesso una penalità adattiva
di Kullback--Leibler verso la policy pre-aggiornamento come seconda
salvaguardia; la configurazione della \cref{sec:optimization} usa
insieme il clip ($\varepsilon = 0.2$) e un target KL di $0.02$.

\paragraph{Generalized Advantage Estimation.}
Le stime di vantaggio $A_t$ provengono da GAE~\cite{schulman2016gae}:
con i residui alle differenze temporali
$\delta_t = r_t + \gamma V(s_{t+1}) - V(s_t)$,
\begin{equation}
  A_t^{\mathrm{GAE}(\gamma,\lambda)} \;=\;
  \sum_{l \geq 0} (\gamma\lambda)^{l}\, \delta_{t+l},
\end{equation}
dove $\lambda \in [0,1]$ scambia bias e varianza tra il TD a un passo
($\lambda = 0$) e i ritorni Monte Carlo ($\lambda = 1$). La piattaforma
usa $\lambda = 0.95$. La qualità di GAE dipende direttamente dalla
qualità del critic --- un episodio di sviluppo documentato
(\cref{sec:iterations}) ha coinvolto un critic la cui varianza spiegata
era $\approx 0$, rendendo i vantaggi quasi privi di baseline finché la
configurazione della value loss non è stata riparata.

\paragraph{Regolarizzazione di entropia.}
Un bonus di entropia aggiunto all'obiettivo sostiene l'esplorazione
penalizzando il collasso prematuro della policy. Il suo coefficiente è
tipicamente ricotto (annealing): alto all'inizio (esplorazione ampia),
basso alla fine (sfruttamento e convergenza stabile). In questa tesi il
coefficiente segue uno schedule esplicito da $0.03$ a $0.005$ ancorato a
una \emph{frazione} fissa del budget di addestramento
(\cref{sec:optimization}), così che run di lunghezza diversa attraversino
lo stesso profilo di esplorazione.

% ------------------------------------------------------------
\section{Apprendimento per rinforzo multi-agente e MAPPO}
\label{sec:marl}

Con $N$ agenti che agiscono in un unico ambiente, l'MDP si generalizza a
un \emph{gioco di Markov}: transizioni e reward dipendono dall'azione
congiunta. Ne seguono due difficoltà strutturali.

\paragraph{Non-stazionarietà.}
Dal punto di vista di ogni singolo agente, gli altri agenti sono parte
dell'ambiente --- e stanno apprendendo, quindi la dinamica effettiva
dell'ambiente cambia nel tempo. Ciò rompe le assunzioni di
stazionarietà alla base delle garanzie di convergenza a singolo agente
e rende fuorviante l'esperienza raccolta sotto vecchie policy dei
co-giocatori~\cite{lowe2017maddpg}. Nella piattaforma del
\cref{ch:system} questo non è un caso limite ma il cuore del setting:
sei agenti che apprendono condividono un mondo, e il \cref{ch:results}
documenta un equilibrio concreto (veicoli che si bloccano in risposta al
comportamento dei pedoni) che esiste solo a causa di questo
accoppiamento.

\paragraph{Attribuzione del merito e varianza.}
Il ritorno di ogni agente dipende dalle azioni di tutti gli agenti,
quindi la varianza delle stime naive di gradiente di policy cresce con
il numero di agenti: un agente può essere premiato o punito per esiti
che non ha causato.

\paragraph{Addestramento centralizzato, esecuzione decentralizzata.}
La risoluzione ormai standard sfrutta la divisione actor--critic:
addestrare i \emph{critic} con accesso a informazione globale (le
osservazioni di tutti gli agenti, o il vero stato), mantenendo gli
\emph{actor} condizionati solo sulle osservazioni
locali~\cite{lowe2017maddpg}. Il critic assorbe la non-stazionarietà ---
condizionato sulla situazione di tutti, il target di valore è meglio
determinato --- mentre l'esecuzione resta pienamente decentralizzata,
poiché i critic vengono scartati dopo l'addestramento. È il paradigma
CTDE usato in tutta questa tesi.

\paragraph{MAPPO.}
MAPPO istanzia il CTDE con actor PPO e una funzione di valore
centralizzata. Yu et al.~\cite{yu2022surprising} hanno mostrato che
questa combinazione concettualmente minimale --- a lungo ritenuta troppo
inefficiente in termini campionari per il MARL --- eguaglia o supera i
metodi off-policy sui benchmark cooperativi standard (particle world,
micromanagement di StarCraft, Hanabi, football), a patto che i dettagli
implementativi siano corretti; tra i fattori che evidenziano ci sono il
design dell'input del critic centralizzato e la normalizzazione del
valore. Questi due fattori ricompaiono concretamente in questa tesi come
il layout dell'osservazione globale della \cref{sec:obs-actions} e
l'episodio di configurazione della value loss della
\cref{sec:iterations}.

\paragraph{Condivisione di parametri e classi eterogenee.}
Agenti omogenei possono condividere una sola rete di policy, mettendo in
comune l'esperienza; agenti eterogenei non possono, perché i loro spazi
di osservazione e azione differiscono. La piattaforma prende la via
intermedia naturale per la sua popolazione: una policy condivisa per
classe (tre veicoli; tre pedoni), due policy in totale
(\cref{sec:architecture}).

\paragraph{Encoder dell'input del critic.}
Il critic centralizzato più semplice appiattisce le osservazioni di
tutti gli agenti in un unico vettore per un MLP. Le alternative
strutturate trattano gli agenti come nodi di un insieme o di un grafo:
gli slot per agente vengono embeddati, combinati per attenzione o per
message passing, e aggregati. La variante su grafo usata come asse
sperimentale secondario di questa tesi applica un'aggregazione a media
in stile GraphSAGE~\cite{hamilton2017graphsage} sui sei nodi-agente con
una maschera di vitalità (\cref{sec:critic-encoders}). Il confronto
critic MLP-vs-GNN chiede se la struttura relazionale nella stima del
valore cambi ciò che gli (identici) actor apprendono.

\paragraph{Coesistenza piuttosto che cooperazione.}
A differenza dei benchmark citati, gli agenti di questa tesi non
condividono una reward di squadra né competono: ognuno persegue il
proprio percorso sotto una reward individuale (\cref{sec:rewards}). Il
CTDE resta lo strumento giusto --- il critic centralizzato stabilizza
comunque la stima del valore sotto interazione reciproca --- ma il
setting è meglio descritto come \emph{coesistenza} di task indipendenti
in un mondo condiviso, l'inquadramento formalizzato come copertura
parallela dei task nella \cref{sec:parallel-coverage}.

% ------------------------------------------------------------
\section{Curriculum learning e la sua condizione di controllo}
\label{sec:curriculum-bg}

\paragraph{Origini.}
L'idea che l'apprendimento migliori quando gli esempi sono presentati in
un ordine sensato dal facile al difficile precede il deep learning: gli
esperimenti ``starting small'' di Elman mostrarono reti ricorrenti che
apprendevano una grammatica complessa solo quando la complessità
dell'input (o la memoria effettiva della rete) cresceva
incrementalmente~\cite{elman1993starting}. Bengio et
al.~\cite{bengio2009curriculum} diedero alla strategia il nome di
\emph{curriculum learning} e la inquadrarono come metodo di
continuazione: gli esempi facili definiscono una versione più liscia
dell'obiettivo su cui l'ottimizzazione parte, quelli difficili vengono
introdotti mano a mano, con guadagni misurabili in velocità di
ottimizzazione e generalizzazione su task supervisionati.

\paragraph{Curriculum learning nel RL.}
Nell'apprendimento per rinforzo il progettista controlla non solo
l'ordine dei dati ma la distribuzione dei \emph{task}, il che rende i
curricula al tempo stesso più potenti e più difficili da progettare.
Narvekar et al.~\cite{narvekar2020curriculum} organizzano il campo
attorno a tre domande: come vengono generati i task sorgente, come
vengono \emph{sequenziati}, e come la conoscenza si trasferisce tra di
essi. In quella tassonomia, il meccanismo della \cref{sec:regimes} è uno
schema di sequenziamento basato sulla prestazione su un insieme fisso di
task: tre livelli di difficoltà ordinati per lunghezza del percorso, con
progressione condizionata a una soglia di successo a finestra e
garantita da cap di pressione basati sul budget. La motivazione dei
curricula nel RL riguarda in larga parte l'esplorazione: sui task
difficili una policy debole può non sperimentare essenzialmente mai il
successo, non ricevendo alcun segnale di apprendimento, mentre i task
facili forniscono un successo iniziale denso da cui la competenza può
avviarsi.

\paragraph{Rischi noti.}
Tre modi di fallimento ricorrono in letteratura e hanno plasmato il
design della \cref{sec:regimes}. Primo, l'\emph{oblio}: le reti neurali
addestrate su una distribuzione spostata tendono a sovrascrivere la
competenza acquisita in precedenza~\cite{french1999catastrophic}; un
curriculum che abbandona i task facili dopo la progressione può erodere
proprio le abilità che ha costruito --- la piattaforma mantiene perciò
tutti i livelli sbloccati nel mix di campionamento invece di
sostituirli. Secondo, lo \emph{stallo}: un gate di competenza mai
soddisfatto può intrappolare l'addestramento sui task facili
indefinitamente --- da qui i cap di pressione che forzano la
progressione entro il budget. Terzo, la \emph{mis-calibrazione del
gate}: sbloccare troppo presto espone la policy a rumore da cui non può
ancora apprendere; sbloccare troppo tardi spreca budget. Questi rischi
sono esattamente le ipotesi che H1 mette alla prova empiricamente
(\cref{sec:rq}).

\paragraph{La condizione di controllo: addestramento misto.}
Ogni affermazione sull'ordinamento richiede una baseline in cui
l'ordinamento è assente. Il controllo naturale è l'addestramento
\emph{misto} (interleaved): tutti i livelli di difficoltà presentati dal
primo episodio, con esposizione uniforme --- l'analogo RL
dell'addestramento su un dataset mescolato, che è il default
nell'apprendimento supervisionato. Il campionatore stratificato della
\cref{sec:batch-arm} implementa esattamente questo controllo
(esposizione uniforme $1/3$ con sbilanciamento massimo di un episodio),
così che il confronto curriculum-contro-misto isoli la variabile di
ordinamento a parità di budget. L'addestramento misto porta con sé la
propria firma ipotizzata (H2): piena copertura dei task fin
dall'inizio, al prezzo di un segnale iniziale più rumoroso dominato dai
fallimenti sui task difficili.

% ------------------------------------------------------------
\section{Il simulatore CARLA}
\label{sec:carla-bg}

CARLA~\cite{dosovitskiy2017carla} è un simulatore open-source di guida
urbana costruito su Unreal Engine~4, sviluppato per supportare
l'addestramento e la validazione di sistemi di guida autonoma. Gira come
server che ospita il mondo 3D e la fisica, controllato da client
attraverso una API Python; è distribuito con un insieme di città di
layout vario, traffico scriptato e meteo configurabile. Tre
caratteristiche contano per questa tesi (versione 0.9.16 ovunque):

\begin{itemize}
  \item \textbf{Modalità sincrona con passo temporale fisso.} Il server
        avanza solo sui tick del client, a passo di fisica fisso
        (\SI{0.05}{s} qui). La raccolta di esperienza è quindi
        riproducibile nella struttura e indipendente dal rendering o
        dalla velocità di orologio (\cref{sec:architecture}).
  \item \textbf{Traffic Manager e controllo dei walker.} I veicoli
        scriptati (NPC) sono guidati dal Traffic Manager di CARLA; i
        pedoni sono attori \emph{walker} che accettano comandi di
        controllo diretti, il che è ciò che rende possibile apprendere
        policy pedonali (\cref{sec:obs-actions}).
  \item \textbf{Città strutturalmente diverse.} L'addestramento usa
        \emph{Town03}, descritta dalla documentazione ufficiale come
        ``a larger, urban map with a roundabout and large
        junctions''~\cite{carladocs}; il test di generalizzazione usa
        \emph{Town05}, una ``squared-grid town with cross junctions and
        a bridge'', con più corsie per direzione~\cite{carladocs}. I
        layout differiscono in natura, non solo in dimensione, rendendo
        il trasferimento Town03$\to$Town05 un test genuino di
        generalizzazione di mappa (\cref{sec:generalization}).
\end{itemize}

Due note di scope. Primo, in questa tesi CARLA fornisce fisica, dinamica
del mondo e traffico --- non percezione: gli agenti osservano feature di
stato privilegiate (\cref{sec:obs-actions}), non flussi camera o lidar,
quindi tutti i risultati riguardano la decisione sotto interazione, non
la visione. Secondo, anche in modalità sincrona e con tutti i seed
fissati, CARLA sotto carico multi-agente non è deterministico da run a
run; questa variabilità irriducibile è quantificata con repliche a seed
identico e usata come metro per ogni effetto riportato in seguito
(\cref{sec:iterations}, \cref{sec:threats}).
